The corrected cup products indexed by two odd primes all lift to full
GT. The construction uses finite homology of cyclic covers and keeps
track of the entire central extension, including its cyclic terms.
For a multiplier-one element \(g\), let \(B_g(X,Y)\) denote the
normalized scalar on the completed Alexander commutator module,
as defined in Section~\ref{sec:all-odd-cross-cup-lifts}.

\begin{theorem}\label{thm:main-odd-pair-lifts}
For every ordered pair of distinct odd primes \(\ell,q\), there is
an invariant continuous homomorphism \(\xi_{\ell,q}:M\to\F_2\),
where \(M=\ker\lambda\), whose restriction to \(N\) belongs to
\(X_2\) and satisfies
\[
                    \xi_{\ell,q}|_K=\kappa_{\ell,q}.
\]
Here \(\kappa_{\ell,q}=\kappa_{q,\ell}\) denotes the basis
character indexed by the unordered pair. Its transgression is
\[
 c_\ell c_q+\ell_\ell(q^*)\delta_\ell
                         +\ell_q(\ell^*)\delta_q.
\]
In particular this corrected class inflates to zero in
\(H^2(G,\F_2)\).

Explicitly, put \(d=\operatorname{ord}_\ell(q)\), and choose a
primitive \(\ell\)-th root \(\zeta\in\F_{q^d}\).
If \(d=2s\), choose representatives \(a_i\) of
\((\Z/\ell)^\times/\langle q\rangle\); then
\[
 (-1)^{\xi_{\ell,q}(g)}
   =\prod_i B_g(\zeta^{a_i},\zeta^{a_i})^{(q^s+1)/2}.
\]
If \(d\) is odd, choose representatives of
\((\Z/\ell)^\times/\langle q,-1\rangle\); then
\[
 (-1)^{\xi_{\ell,q}(g)}
   =\prod_i B_g(\zeta^{a_i},\zeta^{a_i})^{(q^d-1)/2}.
\]
The values are signs independent of these choices.
The same construction works for any continuous positive peripheral
action on \(\widehat F(x,y)\) of the form
\(x\mapsto x^\lambda\), \(y\mapsto f^{-1}y^\lambda f\),
with \(f\in\widehat F(x,y)'\) and
\((xy)^{-1}\mapsto\text{a conjugate of }(xy)^{-\lambda}\).

The two ordered constructions have the same arithmetic restriction.
Their sum on \(N\) extends uniquely to a continuous character
\(\varepsilon_{\ell,q}:G\to\F_2\) vanishing on \(H\).
Whether these relative characters vanish is not determined here.
\end{theorem}

Theorem~\ref{thm:all-odd-cross-cup-lifts} proves the general
construction. The pair \((5,29)\) gives the smallest worked example
used below: its two cyclic correction terms are zero, and the
character detects the classical dihedral field in
Theorem~\ref{thm:main-full-quadratic-target}.
The arithmetic relation between Fermat homology and Anderson's beta
function is classical; Otsubo gives its precise cocycle comparison
in \cite[\S7, p.~637, (7.1)--(7.2)]{Otsubo2016}.
Here the finite symplectic normalizer computes the central extension
for the full peripheral action.
